\section{Targets with small Fourier spectral norm}
\label{sec:fourier-spectral-targets}

For a Boolean function $f:\{-1,1\}^d\to\{-1,1\}$, write
\[
 f(x)=\sum_{S\subseteq[d]}\widehat f(S)\chi_S(x),
 \qquad \chi_S(x)=\prod_{i\in S}x_i,
 \qquad A(f)=\sum_S|\widehat f(S)|.
\]
The Fourier coefficients here use the uniform distribution, even when
the learning distribution is different. Parseval implies $A(f)\ge1$.
The following sufficient condition imposes neither a Fourier-degree
bound nor a bound on the number of nonzero Fourier coefficients.

\begin{lemma}[A shallow realization from Fourier spectral norm]
\label{lem:fourier-spectral-shallow}
Every such $f$, with $A=A(f)$, has a one-hidden-layer tanh realization
with linear readout, signed margin at least $13/8$, width at most
\[
 1+d\left\lceil128A^2(d+2)\log2\right\rceil,
\]
and physical weights bounded by
\[
 \max\{3A,\tfrac d2\log(32Ad),1\}.
\]
\end{lemma}

\begin{proof}
Sample a set $S$ with probability $|\widehat f(S)|/A$ and attach
the sign $s_S=\operatorname{sign}\widehat f(S)$.
The random function $A s_S\chi_S$ has mean $f$ and takes values
in $[-A,A]$. For
$m=\lceil128A^2(d+2)\log2\rceil$, Hoeffding's inequality and a
union bound over the $2^d$ inputs show that, with positive probability,
\[
 \left|\frac A m\sum_{j=1}^m s_{S_j}\chi_{S_j}(x)-f(x)\right|
 \le\frac18\qquad\text{for every }x.
\]
Fix one such choice.

Each parity on $k\ge1$ variables admits an elementary real-output
threshold expansion. Put $t=\sum_{i\in S}x_i$ and
$\theta_j=-k+2j-1$ for $1\le j\le k$. On the permitted values
$t\in\{-k,-k+2,\ldots,k\}$,
\[
 \chi_S(x)=b_k+\sum_{j=1}^k(-1)^{k-j}
                  \operatorname{sign}(t-\theta_j),
 \qquad b_k\in\{0,1\}.
\]
Indeed, both sides start at $(-1)^k$ and have the same jump at
each threshold; explicitly $b_k=1$ for even $k$ and zero for odd $k$.
Every threshold has absolute affine margin at least one.
With $\lambda=\tfrac12\log(32Ad)$, replace each sign by
$\tanh(\lambda(t-\theta_j))$. The resulting parity approximation
has uniform error at most
\[
 2k e^{-2\lambda}\le\frac1{16A}.
\]
The empty parity is the constant one. A single constant hidden neuron,
whose output is $\tanh1$, supplies all constant terms.

Replacing the $m$ sampled parities by these approximations adds error
at most $1/16$. Thus the resulting scalar output is within $3/16$ of
$f$ everywhere. Doubling the readout gives signed margin at least
$13/8$. There are at most $md$ nonconstant hidden neurons. Their
first-layer coefficients and biases have magnitude at most $d\lambda$;
their output coefficients have magnitude at most $2A/m$. The single
constant neuron's output coefficient has magnitude at most
$2A/\tanh1<3A$. These are the stated bounds.
\end{proof}

\begin{corollary}[Unrestricted learning from small spectral norm]
\label{cor:fourier-spectral-unrestricted}
For each fixed $\zeta>0$, there is a polynomial $P_\zeta$ such that
for every Boolean target $f$, every input distribution, and every
learner depth $L\ge2$, a sample size
\[
 n\ge P_\zeta\bigl(d,A(f),L,\varepsilon^{-1},\log(2/\delta)\bigr)
\]
gives, with probability at least $1-\delta$, Gibbs error at most
$\varepsilon$ for the original unrestricted Gaussian posterior,
simultaneously for every $N\ge C_L(f)$ and every
$0<\sigma^2\le N^{-(1+\zeta)}$.
The posterior retains its original conditioning on a global clean
output margin and the training labels.
\end{corollary}

\begin{proof}
Substitute Lemma~\ref{lem:fourier-spectral-shallow} into
Theorem~\ref{thm:shallow-all-depths}. All witness parameters are
polynomial in $d$ and $A(f)$. The all-width conclusion is the
all-width conclusion of that theorem and does not require
$C_L(f)$ to equal the width of the constructed shallow realization.
\end{proof}

Both the realization lemma and its learning corollary also hold with
$A(f)$ replaced by $\sum_S|\widehat g(S)|$ for any real-valued Fourier
polynomial $g$ satisfying $f(x)g(x)\ge1$ on every input. The same
sampling argument approximates $g$ to error $3/16$, which still leaves
the stated signed margin after doubling the readout. Thus a small
normalized coefficient sum for a margin-separated vote of parities
is enough, even when the Boolean sign function itself has a larger
Fourier norm.

For example, every parity has $A(f)=1$, independently of its degree.
More generally, a function with $s$ nonzero Fourier coefficients has
$A(f)\le\sqrt s$ by Cauchy--Schwarz and Parseval. These are sufficient
representation conditions. The argument does not assert that all small
Boolean formulas or all narrow branching programs have polynomial
Fourier spectral norm.

\begin{corollary}[Two-state branching programs]
\label{cor:two-state-branching}
Every Boolean function computed by a deterministic layered width-two
branching program of length $T$ has $A(f)\le T+1$. Consequently the sample bound
in Corollary~\ref{cor:fourier-spectral-unrestricted} is polynomial in
$d,T,L,\varepsilon^{-1},\log(2/\delta)$ for this class.
The program may query variables repeatedly, and its two nodes in a
layer may query different variables.
\end{corollary}

\begin{proof}
Encode the two states as $-1,+1$, and let $s_j(x)$ be the state after
$j$ transitions. Given the old state $+1$ or $-1$, the next state is
a Boolean function $u_+(x)$ or $u_-(x)$ of at most one input bit.
Each such function is constant or a signed input bit, so
$A(u_+)=A(u_-)=1$. The transition can be written
\[
 s_{j+1}(x)=a_j(x)s_j(x)+b_j(x),\qquad
 a_j=(u_+-u_-)/2,\quad b_j=(u_++u_-)/2.
\]
The Fourier coefficient norm is subadditive and submultiplicative:
the latter follows because multiplication of parities uses symmetric
difference of their index sets. Thus $A(a_j),A(b_j)\le1$ and
\[
 A(s_{j+1})\le A(s_j)+1.
\]
The initial state is constant, so $A(s_0)=1$. The final output is
either constant or a signed copy of the final state, proving the claim.
Unused state slots can be filled arbitrarily. If the program uses
absorbing states for early termination, these states count toward its
width-two bound.
\end{proof}

The same proof allows a node to query a parity rather than a single
bit, since a signed parity also has Fourier norm one. Decision lists
with $T$ parity tests likewise have Fourier norm at most $T+1$, by a
separate direct argument: write the output as the sum of its $T+1$
signed first-accept and default indicators. Each indicator is a
product of factors $(1\pm\chi_S)/2$, each of Fourier norm at most
one, so submultiplicativity and the triangle inequality give the
bound. These extensions concern representation and sample bounds
only; no efficient algorithm for identifying a program from samples
is asserted.
